\documentclass[11pt]{article}
\usepackage{mathblatt}
\usepackage{multicol}


\begin{document}
\pfheftstil
\pfheftkopf{Prozent \textperiodcentered{} Lösungen}{}
\begin{pfloesung}{}
\lz{1.}{50/100, 25/100, 20/100, 75/100, $\tfrac{30}{100}$}{}{}
\end{pfloesung}
\begin{pfloesung}{}
\lz{2.}{1 \% = $\tfrac{1}{100}$ = 0,01; 10 \% = $\tfrac{1}{10}$ = 0,1; 20 \% = $\tfrac{1}{5}$ = 0,2; 25 \% = $\tfrac{1}{4}$ = 0,25; 50 \% = $\tfrac{1}{2}$ = 0,5; 100 \% = 1 = 1}{}{}
\end{pfloesung}
\begin{pfloesung}{}
\lz{3.}{7; 60; 10}{}{}
\end{pfloesung}
\begin{pfloesung}{}
\lz{4.}{$\tfrac{1}{8}$ = 0,125 = 12,5 \%; $\tfrac{3}{5}$ = 0,6 = 60 \%; $\tfrac{7}{20}$ = 0,35 = 35 \%; $\tfrac{2}{3}$ = 0,666… $\approx$ 66,7 \%}{}{}
\end{pfloesung}
\begin{pfloesung}{}
\lz{5.}{7 €}{}{}
\end{pfloesung}
\begin{pfloesung}{}
\lz{6.}{+20 \% $\to$ 1,2; $-$20 \% $\to$ 0,8; +8 \% $\to$ 1,08; $-$11 \% $\to$ 0,89; +1,9 \% $\to$ 1,019}{}{}
\end{pfloesung}
\begin{pfloesung}{}
\lz{7.}{$55\,\%$}{$\tfrac{55}{100}$}{}
\end{pfloesung}
\begin{pfloesung}{}
\lz{8.}{jeder $5.$}{$20\,\% = \tfrac{20}{100} = \tfrac{1}{5}$}{}
\end{pfloesung}
\begin{pfloesung}{}
\lz{9.}{15 \% = $\tfrac{15}{100}$ = $\tfrac{3}{20}$; 0,7 = $\tfrac{7}{10}$}{}{}
\end{pfloesung}
\begin{pfloesung}{}
\lz{10.}{$\tfrac{14}{40}$ und 35 \%}{}{}
\end{pfloesung}
\begin{pfloesung}{}
\lz{11.}{4 von 100 Bäumen wachsen nicht an}{}{}
\end{pfloesung}
\begin{pfloesung}{}
\lz{12.}{$50\,\%$}{}{}
\end{pfloesung}
\begin{pfloesung}{}
\lz{13.}{20 \%}{$\tfrac{1}{5}$ = $\tfrac{20}{100}$}{}
\end{pfloesung}
\begin{pfloesung}{}
\lz{14.}{$10$ gleiche Teile}{}{}
\end{pfloesung}
\begin{pfloesung}{}
\lz{15.}{$28$ Kinder}{}{}
\end{pfloesung}
\begin{pfloesung}{}
\lz{16.}{20 000}{}{}
\end{pfloesung}
\begin{pfloesung}{}
\lz{17.}{$12\,€$}{ein Viertel vom Ganzen: $48 : 4$}{}
\end{pfloesung}
\begin{pfloesung}{}
\lz{18.}{20 \% von 50 € = 10 €}{}{}
\end{pfloesung}
\begin{pfloesung}{}
\lz{19.}{110 €}{}{}
\end{pfloesung}
\begin{pfloesung}{}
\lz{20.}{90 €}{}{}
\end{pfloesung}
\begin{pfloesung}{}
\lz{21.}{24 €}{}{}
\end{pfloesung}
\begin{pfloesung}{}
\lz{22.}{21 €}{W = G · p = 70 € · 0,3}{}
\end{pfloesung}
\begin{pfloesung}{}
\lz{23.}{$48\,€$}{Rabatt $64 : 4 = 16\,€$; neuer Preis $64 - 16$}{}
\end{pfloesung}
\begin{pfloesung}{}
\lz{24.}{2 500 · 0,1 = 250}{}{}
\end{pfloesung}
\begin{pfloesung}{}
\lz{25.}{190 €}{}{}
\end{pfloesung}
\begin{pfloesung}{}
\lz{26.}{5 \% von 600 € = 30 €}{}{}
\end{pfloesung}
\begin{pfloesung}{}
\lz{27.}{12 \% von 75 = 9}{}{}
\end{pfloesung}
\begin{pfloesung}{}
\lz{28.}{$2{,}80\,€$}{$0{,}07 \cdot 40$}{}
\end{pfloesung}
\begin{pfloesung}{}
\lz{29.}{3 785,40 € (Rabatt 420,60 €)}{}{}
\end{pfloesung}
\begin{pfloesung}{}
\lz{30.}{67,20 €}{}{}
\end{pfloesung}
\begin{pfloesung}{}
\lz{31.}{$0{,}90\,€$}{$0{,}2 \cdot 4{,}50$}{}
\end{pfloesung}
\begin{pfloesung}{}
\lz{32.}{Term 1 falsch; Term 2 richtig; Term 3 richtig}{voller Preis 3 · 12 € + 7,50 € = 43,50 €; Rabatt 8,70 €}{}
\end{pfloesung}
\begin{pfloesung}{}
\lz{33.}{0,7 · 49 cm³ = 34,3 cm³}{}{}
\end{pfloesung}
\begin{pfloesung}{}
\lz{34.}{24 \% von 700 = 168}{}{}
\end{pfloesung}
\begin{pfloesung}{}
\lz{35.}{14 560 € · 0,16 = 2 329,60 €}{}{}
\end{pfloesung}
\begin{pfloesung}{}
\lz{36.}{6,50 €}{}{}
\end{pfloesung}
\begin{pfloesung}{}
\lz{37.}{15 \% von 440 cm³ = 66 cm³}{}{}
\end{pfloesung}
\begin{pfloesung}{}
\lz{38.}{4 331,60 €}{}{}
\end{pfloesung}
\begin{pfloesung}{}
\lz{39.}{208 000 · 0,141 = 29 328}{}{}
\end{pfloesung}
\begin{pfloesung}{}
\lz{40.}{276 Jugendliche}{23 \% von 1200 = 276}{}
\end{pfloesung}
\begin{pfloesung}{}
\lz{41.}{$20\,\%$}{}{}
\end{pfloesung}
\begin{pfloesung}{}
\lz{42.}{$17\,\%$}{Hundertstel: $\tfrac{17}{100}$}{}
\end{pfloesung}
\begin{pfloesung}{}
\lz{43.}{30 \%}{}{}
\end{pfloesung}
\begin{pfloesung}{}
\lz{44.}{30 \%}{}{}
\end{pfloesung}
\begin{pfloesung}{}
\lz{45.}{$30\,\%$}{Ganzes $20$; $6 : 20 = 0{,}3$}{}
\end{pfloesung}
\begin{pfloesung}{}
\lz{46.}{$25\,\%$}{Rabatt $20\,€$; $20 : 80 = 0{,}25$}{}
\end{pfloesung}
\begin{pfloesung}{}
\lz{47.}{12,5 \%}{}{}
\end{pfloesung}
\begin{pfloesung}{}
\lz{48.}{am Dienstag}{Montag $36 : 240 = 0{,}15 = 15\,\%$; Dienstag $27 : 150 = 0{,}18 = 18\,\%$}{}
\end{pfloesung}
\begin{pfloesung}{}
\lz{49.}{70 : 400 = 0,175 = 17,5 \%}{}{}
\end{pfloesung}
\begin{pfloesung}{}
\lz{50.}{$\approx$ 14,4 \%}{}{}
\end{pfloesung}
\begin{pfloesung}{}
\lz{51.}{4 613 : 15 000 $\approx$ 30,75 \%}{}{}
\end{pfloesung}
\begin{pfloesung}{}
\lz{52.}{9,96 € : 830 € = 0,012 = 1,2 \%}{}{}
\end{pfloesung}
\begin{pfloesung}{}
\lz{53.}{5 575 : 8 200 $\approx$ 0,68 = 68 \%}{}{}
\end{pfloesung}
\begin{pfloesung}{}
\lz{54.}{165 : 200 = 0,825 = 82,5 \%}{}{}
\end{pfloesung}
\begin{pfloesung}{}
\lz{55.}{$\tfrac{29}{72}$ $\approx$ 40,3 \%}{145 : 360 = $\tfrac{29}{72}$ $\approx$ 0,4028}{}
\end{pfloesung}
\begin{pfloesung}{}
\lz{56.}{$\approx$ 45,2 \%}{}{}
\end{pfloesung}
\begin{pfloesung}{}
\lz{57.}{$\approx$ 26,2 \%}{}{}
\end{pfloesung}
\begin{pfloesung}{}
\lz{58.}{37,7 : 144 $\approx$ 0,262 = 26,2 \%}{}{}
\end{pfloesung}
\begin{pfloesung}{}
\lz{59.}{$\tfrac{17}{24} \approx 70{,}8\,\%$}{$17 : 24 = 0{,}7083\ldots$}{}
\end{pfloesung}
\begin{pfloesung}{}
\lz{60.}{$\approx$ 81,9 \%}{}{}
\end{pfloesung}
\begin{pfloesung}{}
\lz{61.}{(117,92 $-$ 42,79) : 117,92 $\approx$ 63,71 \%}{}{}
\end{pfloesung}
\begin{pfloesung}{}
\lz{62.}{$\tfrac{12}{83}$ $\approx$ 14,5 \%; Sektor $\approx$ 52°}{83 $-$ 71 = 12; 12 : 83 = $\tfrac{12}{83}$ $\approx$ 0,145; $\tfrac{12}{83}$ · 360° = $\tfrac{4320}{83}$° $\approx$ 52°}{}
\end{pfloesung}
\begin{pfloesung}{}
\lz{63.}{$800\,€$}{$1\,\% = 8\,€$; $100\,\%$}{}
\end{pfloesung}
\begin{pfloesung}{}
\lz{64.}{$900\,€$, $450\,€$, $180\,€$}{Ganzes bei $10\,\%$: $90 \cdot 100 : 10 = 900$; Ganzes bei $20\,\%$: $90 \cdot 100 : 20 = 450$; Ganzes bei $50\,\%$: $90 \cdot 100 : 50 = 180$}{}
\end{pfloesung}
\begin{pfloesung}{}
\lz{65.}{$240$ kg}{wie oft bis hundert Prozent: $100 : 10 = 10$; mal nehmen: $10 \cdot 24 = 240$ kg}{}
\end{pfloesung}
\begin{pfloesung}{}
\lz{66.}{440 € : 1,1 = 400 €}{}{}
\end{pfloesung}
\begin{pfloesung}{}
\lz{67.}{3 000 € : 1,2 = 2 500 €}{}{}
\end{pfloesung}
\begin{pfloesung}{}
\lz{68.}{$60$ l}{$5 \cdot 12$}{}
\end{pfloesung}
\begin{pfloesung}{}
\lz{69.}{800 €}{G = W : p = 200 € : 0,25}{}
\end{pfloesung}
\begin{pfloesung}{}
\lz{70.}{48 € : 0,8 = 60 €}{}{}
\end{pfloesung}
\begin{pfloesung}{}
\lz{71.}{G = 45 : 0,6 = 75}{}{}
\end{pfloesung}
\begin{pfloesung}{}
\lz{72.}{40 \% = 120 $\Rightarrow$ 10 \% = 30 $\Rightarrow$ 100 \% = 300 Kinder}{}{}
\end{pfloesung}
\begin{pfloesung}{}
\lz{73.}{30 €}{G = W : p = 6 € : 0,2}{}
\end{pfloesung}
\begin{pfloesung}{}
\lz{74.}{auf $80\,\%$ gesenkt}{}{}
\end{pfloesung}
\begin{pfloesung}{}
\lz{75.}{$120$ Mitglieder}{steht hinten}{}
\end{pfloesung}
\begin{pfloesung}{}
\lz{76.}{$77\,€$}{Prozentwert: $10\,\% = 7\,€$; dazu: $70 + 7 = 77\,€$}{}
\end{pfloesung}
\begin{pfloesung}{}
\lz{77.}{$60\,€$}{$66 : 1{,}1$}{}
\end{pfloesung}
\begin{pfloesung}{}
\lz{78.}{$25\,\%$}{$10 : 40 = 0{,}25$}{}
\end{pfloesung}
\begin{pfloesung}{}
\lz{79.}{120 € · 0,75 = 90 €}{}{}
\end{pfloesung}
\begin{pfloesung}{}
\lz{80.}{15 · 1,2 = 18 Platten}{}{}
\end{pfloesung}
\begin{pfloesung}{}
\lz{81.}{$50\,\%$ mehr}{$1$ Prozentpunkt; $1 : 2 = 0{,}5$}{}
\end{pfloesung}
\begin{pfloesung}{}
\lz{82.}{670 000 · 1,5 = 1 005 000}{}{}
\end{pfloesung}
\begin{pfloesung}{}
\lz{83.}{156 \%}{}{}
\end{pfloesung}
\begin{pfloesung}{}
\lz{84.}{$15\,\%$}{Veränderung gesucht: $(4{,}83 - 4{,}20) : 4{,}20 = 0{,}15$}{}
\end{pfloesung}
\begin{pfloesung}{}
\lz{85.}{25 \%}{400 000 $-$ 320 000 = 80 000}{}
\end{pfloesung}
\begin{pfloesung}{}
\lz{86.}{4,20 €}{20 \% von 3,50 € = 0,70 €}{}
\end{pfloesung}
\begin{pfloesung}{}
\lz{87.}{19 900 € · 0,88 = 17 512 €}{}{}
\end{pfloesung}
\begin{pfloesung}{}
\lz{88.}{$50\,€$}{$53{,}50 : 1{,}07$}{}
\end{pfloesung}
\begin{pfloesung}{}
\lz{89.}{Ja}{erster Rabatt: $600 \cdot 0{,}9 = 540\,€$; zweiter Rabatt: $540 \cdot 0{,}9 = 486\,€$; vergleichen: $486\,€ < 490\,€$; Rest: $490 - 486 = 4\,€$}{}
\end{pfloesung}
\begin{pfloesung}{}
\lz{90.}{280 $-$ 130 = 150; 210 + 210 : 3 = 280}{ein Drittel von 210 = 70}{}
\end{pfloesung}
\begin{pfloesung}{}
\lz{91.}{34,80 €}{43,50 € · 0,8 = 34,80 €}{}
\end{pfloesung}
\begin{pfloesung}{}
\lz{92.}{$\approx$ 35,5 \%}{}{}
\end{pfloesung}
\begin{pfloesung}{}
\lz{93.}{$\approx$ 32,4 \%}{}{}
\end{pfloesung}
\begin{pfloesung}{}
\lz{94.}{$\approx$ 31,3 \%}{}{}
\end{pfloesung}
\begin{pfloesung}{}
\lz{95.}{$\approx$ 134 \%}{}{}
\end{pfloesung}
\begin{pfloesung}{}
\lz{96.}{$\approx$ 4,1 \%}{}{}
\end{pfloesung}
\begin{pfloesung}{}
\lz{97.}{$\tfrac{5}{179}$ $\approx$ 2,8 \%}{Differenz 0,05 €; 0,05 : 1,79 = $\tfrac{5}{179}$ $\approx$ 0,0279}{}
\end{pfloesung}
\begin{pfloesung}{}
\lz{98.}{$\tfrac{7}{47}$ $\approx$ 14,9 \%}{Differenz 70; 70 : 470 = $\tfrac{7}{47}$ $\approx$ 0,149}{}
\end{pfloesung}
\begin{pfloesung}{}
\lz{99.}{99 cm³ · 0,2 = 19,8 cm³}{}{}
\end{pfloesung}
\begin{pfloesung}{}
\lz{100.}{2,7 Mio.; 8,7 Mio.; $\tfrac{11}{79}$ $\approx$ 13,9 \%}{9,7 $-$ 7,0 = 2,7; Summe 78,3 Mio.; 78,3 : 9 = 8,7; 9,0 $-$ 7,9 = 1,1 Mio.}{}
\end{pfloesung}
\begin{pfloesung}{}
\lz{101.}{nein}{Studie 1 54 · 1,03$^{14}$ $\approx$ 81,7 gegen Studie 2 = 75,6 (Tausend Liter); 54 · 1,03$^{14}$ $\approx$ 81,7; 54 · 1,4 = 75,6}{}
\end{pfloesung}
\begin{pfloesung}{}
\lz{102.}{1~$p$, 2~$W$, 3~$G$, 4~neuer Wert, 5~$p$, 6~$W$}{}{}
\end{pfloesung}
\begin{pfloesung}{}
\lz{103.}{$7\,\%$}{$14 : 200 = 0{,}07$}{}
\end{pfloesung}
\begin{pfloesung}{}
\lz{104.}{in der 7a ist der Anteil größer}{7a: $\tfrac{18}{25} = 72\,\%$; 7b: $\tfrac{21}{30} = \tfrac{7}{10} = 70\,\%$}{}
\end{pfloesung}
\begin{pfloesung}{}
\lz{105.}{Die Gesamtzahl der Teilnehmer im Vorjahr ist unbekannt; 22 \% einer größeren Zahl können mehr sein als 24 \% von 150.}{}{}
\end{pfloesung}
\begin{pfloesung}{}
\lz{106.}{Nein}{Es gibt viel mehr junge Westdeutsche als Ostdeutsche; die Anteile müssen gewichtet werden (Ergebnis liegt näher an 69 \%).}{}
\end{pfloesung}
\begin{pfloesung}{}
\lz{107.}{nein}{120 · 0,7 · 0,8 = 67,20 € > 60 €}{}
\end{pfloesung}
\begin{pfloesung}{}
\lz{108.}{Aussage 1 falsch; etwa vier Fünftel (81 \%) der Mädchen nutzen einen Fernseher}{drei Viertel = 75 \%}{}
\end{pfloesung}
\begin{pfloesung}{}
\lz{109.}{100 m und 6 m; Ja, Steigung 16 : 170 = $\tfrac{8}{85}$ $\approx$ 9,4 \% > 6 \%}{16 : 170 = $\tfrac{8}{85}$ $\approx$ 0,0941}{}
\end{pfloesung}
\begin{pfloesung}{}
\lz{110.}{34,4 \% + 11,2 \% + 5,0 \% = 50,6 \% > 50 \%}{}{}
\end{pfloesung}
\begin{pfloesung}{}
\lz{111.}{A = $\pi$ · 4² · 70 : 360 $\approx$ 9,77 cm²; Anteil $\approx$ 22,43 \%}{}{}
\end{pfloesung}
\begin{pfloesung}{}
\lz{112.}{Ja}{12 Deckel, Abfall 800 $-$ 192$\pi$ $\approx$ 196,8 cm² $\approx$ 24,6 \% $\le$ 25 \%; 100 : 8 = 12,5 $\Rightarrow$ 12 Deckel; Deckelfläche 12 · $\pi$ · 16 = 192$\pi$ $\approx$ 603,2 cm²; Blech 800 cm²}{}
\end{pfloesung}
\begin{pfloesung}{}
\lz{113.}{die Zinsen}{Prozentwert}{}
\end{pfloesung}
\begin{pfloesung}{}
\lz{114.}{2,3 \%}{}{}
\end{pfloesung}
\begin{pfloesung}{}
\lz{115.}{200 €}{}{}
\end{pfloesung}
\begin{pfloesung}{}
\lz{116.}{700\,€}{1\,\% = 7\,€}{}
\end{pfloesung}
\begin{pfloesung}{}
\lz{117.}{8 €}{}{}
\end{pfloesung}
\begin{pfloesung}{}
\lz{118.}{stimmt}{$300 \cdot 0{,}03 = 9{,}00$\,€; oder $9 : 300 = 0{,}03 = 3\,\%$}{}
\end{pfloesung}
\begin{pfloesung}{}
\lz{119.}{1,25 \%}{}{}
\end{pfloesung}
\begin{pfloesung}{}
\lz{120.}{11 000 · 0,0144 : 12 = 13,20 €}{}{}
\end{pfloesung}
\begin{pfloesung}{}
\lz{121.}{200 € · 0,02 = 4,00 €}{}{}
\end{pfloesung}
\begin{pfloesung}{}
\lz{122.}{vom neuen Guthaben 612\,€}{}{}
\end{pfloesung}
\begin{pfloesung}{}
\lz{123.}{5\,050\,€}{Zinsen 50\,€}{}
\end{pfloesung}
\begin{pfloesung}{}
\lz{124.}{824,32 €; 16,49 €}{}{}
\end{pfloesung}
\begin{pfloesung}{}
\lz{125.}{Guthaben = 477,00\,€, Zinsen = 9,54\,€}{}{}
\end{pfloesung}
\begin{pfloesung}{}
\lz{126.}{1\,003,97\,€}{Guthaben: 732,00\,€; 866,64\,€}{}
\end{pfloesung}
\begin{pfloesung}{}
\lz{127.}{Ja, knapp: nach 4 Jahren $\approx$ 2\,251,02\,€}{}{}
\end{pfloesung}
\begin{pfloesung}{}
\lz{128.}{1 850 · 1,009³ $\approx$ 1 900,40 €}{}{}
\end{pfloesung}
\begin{pfloesung}{}
\lz{129.}{$\approx$ 64 022,99 €}{}{}
\end{pfloesung}
\begin{pfloesung}{}
\lz{130.}{$\approx$ 477,66 €}{}{}
\end{pfloesung}
\begin{pfloesung}{}
\lz{131.}{614 · 1,1225$^{5}$ $\approx$ 1 094 Mio. €}{}{}
\end{pfloesung}
\begin{pfloesung}{}
\lz{132.}{5 000 · 1,055³ $-$ 5 000 $\approx$ 871,21 €}{}{}
\end{pfloesung}
\begin{pfloesung}{}
\lz{133.}{$\approx$ 15 685,18 €}{}{}
\end{pfloesung}
\begin{pfloesung}{}
\lz{134.}{$\approx$ 9 193,26 €}{}{}
\end{pfloesung}
\begin{pfloesung}{}
\lz{135.}{1000 € · 1,03$^{5}$ $\approx$ 1159,27 €}{Wachstumsfaktor 1,03}{}
\end{pfloesung}
\begin{pfloesung}{}
\lz{136.}{1~$Z$, 2~Zinseszins, 3~$p$, 4~$Z$, 5~Zinseszins, 6~$p$}{}{}
\end{pfloesung}
\begin{pfloesung}{\textbf{Prüfstein}}
\lz{a)}{$\tfrac{25}{156}$ $\approx$ 16,0 \%}{1,25 : 7,8 = $\tfrac{25}{156}$ $\approx$ 0,160}{}
\lz{b)}{Aussage 2 falsch; Ein Viertel der Weltbevölkerung waren 2020 Kinder (1,95 : 7,8 = 0,25)}{mittleres Alter 3,9 : 7,8 = 50 \%; Kinder 25 \%}{}
\lz{c)}{Minimum 290 Mio. (Arabisch); Maximum 1300 Mio. (Chinesisch); Spannweite 1010 Mio.}{1300 $-$ 290}{}
\lz{d)}{1 Kästchen = 50 Mio. (z. B. 100, 200, … je zwei Kästchen); Hindi; Säule Arabisch 290 : 50 = 5,8 Kästchen hoch}{389 Mio. $\mathrel{\widehat{=}}$ 7,8 Kästchen $\Rightarrow$ 50 Mio. je Kästchen; 10,5 · 50 = 525}{}
\end{pfloesung}
\end{document}